\section{Introduction}
\label{sec:intro}

Information encoded with reference to the self is remembered better than
information processed for meaning, sound, or surface form. This pattern,
the \emph{Self-Reference Effect} (SRE), has been documented since
Rogers, Kuiper, and Kirker \cite{rogers1977sre} and confirmed in the
meta-analysis of Symons and Johnson \cite{symons1997meta} with a robust
Cohen's $d \approx 0.50$ separation between self-referenced and
semantically encoded items. The self-memory system has been a productive framework for connecting
the effect to autobiographical memory and the broader self-concept
\cite{conway2005smemsys}.

What computational mechanism produces it?
Existing models stop short. Retrieved-context theories of episodic
memory --- TCM \cite{howard2002tcm}, CMR \cite{polyn2009cmr}, and the
emotional extension CMR3 \cite{cohen2022cmr3} --- model emotion-modulated
recall but lack a representation of self. Activation-based theories like
ACT-R \cite{anderson2007actr} treat all chunks symmetrically.
Recent cognitive-architecture proposals for LLM agents
\cite{sumers2024cogarch} explicitly enumerate working, episodic,
semantic, and procedural stores; concrete implementations
\cite{memos2025} retrieve via embedding similarity plus heuristic
decay, but none distinguish self-relevance from topical relevance. Conversely, the Free Energy Principle (FEP) literature on
self \cite{friston2010fep,apps2014freeself,limanowski2013minimalself}
articulates a generative model of selfhood but does not connect it to
memory dynamics: precision-weighted self priors have been linked to body
recognition and minimal phenomenal selfhood, not to encoding depth or
forgetting rate.

The white space is the intersection. We propose to model the self as
a generative model $\Mself = (\muself, \piself)$ under the FEP and to let
its precision multiplicatively gate memory:

\begin{equation}
\text{forget\_score}\;=\;
\underbrace{\text{recency}\cdot\text{usage}}_{\text{Ebbinghaus}}
\;\cdot\;\frac{1}{1+\text{importance}}
\;\cdot\;\frac{1}{1+\alpha|\valence|\arousal}
\;\cdot\;\frac{1}{1+\lambda\,\srel}
\end{equation}

where $\srel = \piself \cdot \tfrac{1+\cos(e_i,\muself)}{2}
+ (1-\piself)\cdot 0.5$ is the self-relevance of memory $i$ given its
embedding $e_i$.
The self prior $\muself$ is updated by exponential moving average toward
\emph{identity-constitutive} utterances --- user content containing
first-person tokens --- gated by emotional significance. The precision
$\piself$ follows an FEP-inspired rule
$\piself \propto 1/(1+\gamma\,\Var[\PE])$ over a rolling prediction-error
window.

The contributions of this paper are:

\begin{enumerate}[leftmargin=*,itemsep=2pt,topsep=2pt]
\item \textbf{A precision-weighted generative self that drives memory
dynamics (\cref{sec:method}).} The self enters the existing
forget\_score formula as a multiplicative factor, dual to the emotional
intensity factor; entering it into retrieval ranking adds a single
weighted similarity term.
\item \textbf{Controlled self-relevance response (E1, \cref{sec:exp}).}
On a $4$-condition encoding paradigm with condition-typed self-relevance
values drawn from depth-of-processing theory, the forget-score ordering
follows the human SRE direction in $10/10$ seeds.
We are explicit that $\srel$ is manually assigned, so this tests the
gating formula's qualitative response --- not the stimulus-to-memory
forward pass.
\item \textbf{Non-parallel raw-scale cell means (E2, \cref{sec:exp}).}
A $2{\times}2$ factorial recovers the multiplicative structure analytically
expected from \cref{eq:forget}: cell-mean differentials are not constant
(diagnostic magnitude $0.71$), so the raw-scale linear-additive model is
rejected against the model containing the interaction term. We present the
analytic interpretation alongside the AIC value, because the AIC magnitude
itself is fragile under deterministic data.
\item \textbf{Monotone $\piself$ attenuation (E3, \cref{sec:exp}).} A
true precision ablation through the \texttt{SelfModel} forward pass
yields SRE size $0.000$ at $\piself{=}0$, growing monotonically to
$0.085$ at $\piself{=}1$. This is consistent with empirical SRE
attenuation reported in dissociation and severe depression.
\item \textbf{Open-source substrate.} The complete implementation is a
$\sim$200-line extension to an existing open-source cognitive-layer
substrate \cite{anon_substrate}; all experiments are CPU-only and run
in seconds.
\end{enumerate}

\Cref{sec:limitations} states five scope boundaries together with the
validation roadmap they imply: a bag-of-tokens embedding that defers a
true raw-stimulus SRE forward pass (L1) to follow-up with a
sentence-transformer encoder; a precision-inspired gating rule that is
not yet derived from a variational objective (L2); the absence of a
quantitative fit to human meta-analytic effect sizes (L3). The present
work establishes the mechanism and its falsifiable predictions; the
empirical quantitative fit is the explicit next step.
